\begin{abstract}
Accurate real-time thermal state estimation in lithium-ion battery cells is essential for the safety and efficiency of electric vehicle (EV) battery management systems. Traditional thermal modeling relies either on computationally intensive numerical methods unsuitable for real-time execution, or on simplified lumped-parameter models that struggle under dynamic transients. While data-driven machine learning models offer rapid evaluation, they are vulnerable to exposure bias, where single-step prediction errors compound during autoregressive multi-step rollouts. This paper presents a Physics-Informed Neural Network (PINN) framework for dynamic battery cell temperature prediction that addresses these challenges. This work integrates physics-informed regularization, exposure-bias mitigation through input noise injection, trainable physical parameter estimation, and multi-step rollout evaluation into a unified battery thermal digital twin. Evaluated on a dataset comprising 522 independent dynamic simulations across multiple drive cycles, ambient temperatures (0$^\circ$C--50$^\circ$C), and mass scales (totaling 78 million time steps), the proposed model reduces the maximum rollout error on an unseen test drive cycle (WLTP2) generated within the same Simulink framework from 0.2474$^\circ$C in the baseline MLP to 0.0851$^\circ$C, achieving an autoregressive rollout Mean Absolute Error (MAE) of 0.0117$^\circ$C. The trainable convective parameter ($hA$) converges to an effective value of 36.49~W/K, consistent with the assumed single-node thermal model and surface cooling dynamics.
\end{abstract}

\begin{IEEEkeywords}
Battery thermal digital twin, physics-informed neural networks, exposure bias, parameter estimation, lithium-ion battery, electric vehicles.
\end{IEEEkeywords}
